\section{Areal Law of the Triadic Cut: Capacity \(81\log 3\), Area Operator, Spectrum, and Collective Intertwiner of the Channels \(90/120\)}
\label{hap:sec:area-entrelazamiento-barbero-rev11}
\index{Barbero--Immirzi!area operator}
\index{modular entropy!modular Hamiltonian}

The HMT identification of the functional \(\Gamma_{6\to8}\) with the
Barbero--Immirzi coordinate is made on the same triadic boundary that supports the
incidence (90/120). The interior and its boundary sheet have equal cardinality,
\[
 \mathcal B=(\mathbb Z/9\mathbb Z)^3,
 \quad |\mathcal B|=729,
 \qquad
 \partial\mathcal B=(\mathbb Z/27\mathbb Z)^2,
 \quad |\partial\mathcal B|=729,
\]
by regrouping six trits into three nonadic coordinates or two
coordinates of order twenty-seven.

\begin{theorem}[Areal Law of Triadic Cutoff]
\label{hap:thm:ley-areal-corte-triadico-rev11}
In the cubic graph \(N\times N\times N\) with unit capacity, the minimum
cut between two opposite faces has capacity \(N^2\). For \(N=9\),
\[
 \mathcal A_{\min}=81,
 \qquad S_{\rm tri}=\mathcal A_{\min}\log3=81\log3.
\]
\end{theorem}
\begin{proof}
The \(N^2\) lines parallel to the normal direction are edge-disjoint
paths, so every cut intercepts them. A transverse layer contains
exactly \(N^2\) edges and attains the bound. Each triadic link contributes
\(\log3\) to the entropy of a maximally mixed state.
\end{proof}

On the boundary torus \(27\times27\), let \(A\) be the canonical block
\(9\times9\). Its boundary decomposes exactly into
\begin{equation}
 \partial A=\partial A_{90}\sqcup\partial A_{120},
 \qquad |\partial A_{90}|=|\partial A_{120}|=18.
 \label{hap:eq:corte-90-120-rev11}
\end{equation}
The horizontal links represent channel (90), and the vertical ones
channel (120). This decomposition turns the pair of incidences into an
operator decomposition of the boundary.

The areal entropy \(S_{\rm tri}=81\log3\) belongs to the maximally
mixed state of the cut links. The bilayer entropy
\(S_{\rm bi}(y)\) of the \cref{hap:def:entropia-bicapa-rev6} belongs to the
normalized distribution between the two sheets; the two objects retain
distinct domains.

Let \(N=\operatorname{diag}(0,1,2)\) be the number operator of a tritic link.
The areal realization declares the positive scale
\[
 \theta_{\rm clk}=\frac{C^*}{6}\frac{\pi}{180},
 \qquad
 a_0=\frac{1+2\cos(10^\circ)}{3}\,\theta_{\rm clk},
 \qquad \gamma_{\rm HMT}=\Gamma_{6\to8}.
\]
The formula for \(a_0\) constitutes the starting structure of this
realization; once that scale is fixed, the operator and its spectrum follow
exactly.

\begin{definition}[Area Operator HMT]
On the boundary sheet
\(\mathcal H_{\partial}=\bigotimes_{e\in\partial A}\mathbb C^3\), define
\begin{equation}
 \widehat{\mathcal A}_{\rm HMT}
 =\gamma_{\rm HMT}a_0\sum_{e\in\partial A}N_e.
 \label{hap:eq:operador-area-hmt-rev11}
\end{equation}
Its spectrum in the canonical cut is
\[
 \operatorname{Spec}\widehat{\mathcal A}_{\rm HMT}
 =\{n\gamma_{\rm HMT}a_0:0\leq n\leq72\},
\]
with multiplicities given by the coefficients of
\((1+z+z^2)^{36}\). Prolongation through compatible cuts produces the
protected family \(\mathcal A_n^{\rm prot}=n\gamma_{\rm HMT}a_0\).
\end{definition}

The intertwiner that fixes normalization is constructed in the subspace of
collective channels. Let \(u_{90},u_{120}\) be the uniform unit vectors
of the hexadic and octadic incidences, and let \(v_{90},v_{120}\) be the
uniform unit vectors of the two sets of links of
\eqref{hap:eq:corte-90-120-rev11}. The map
\begin{equation}
 \begin{aligned}
 \mathcal J_{6\to8}:{}
 &\operatorname{span}\{u_{90},u_{120}\}
 \longrightarrow
 \operatorname{span}\{v_{90},v_{120}\},\\
 &\mathcal J_{6\to8}u_{90}=v_{90},
 \qquad
 \mathcal J_{6\to8}u_{120}=v_{120}
 \end{aligned}
 \label{hap:eq:entrelazador-area-barbero-rev11}
\end{equation}
is an isometry between these collective subspaces. If
\[
 D_{\rm inc}=90|u_{90}\rangle\langle u_{90}|
 +120|u_{120}\rangle\langle u_{120}|,
 \qquad
 D_{\rm area}=90|v_{90}\rangle\langle v_{90}|
 +120|v_{120}\rangle\langle v_{120}|,
\]
then
\[
 \boxed{\mathcal J_{6\to8}D_{\rm inc}
 =D_{\rm area}\mathcal J_{6\to8}.}
\]
In particular, the primitive combination \(12:-1\) acts
with the same coefficients before and after \(\mathcal J_{6\to8}\). This
equality fixes the normalization dictionary that inserts
\(\gamma_{\rm HMT}=\Gamma_{6\to8}\) into the
Ashtekar--Barbero connection.

